\documentclass{article}

\usepackage[utf8]{inputenc}
\usepackage[T1]{fontenc}
\usepackage{amsmath,amssymb,amsthm}
\usepackage{mathtools}
\usepackage{hyperref}
\usepackage{booktabs}

\makeatletter
\newcommand{\citep}[2][]{%
  \begingroup
  \def\@tmpkey{#2}%
  \ifx\@tmpkey\@empty\else
    (\ifx\@empty#1\@empty\else #1, \fi
     \@nameuse{bibkey@\@tmpkey})%
  \fi
  \endgroup
}
\expandafter\def\csname bibkey@lasser2026exo\endcsname{Exogenous~Verification}
\expandafter\def\csname bibkey@lasser2026phase\endcsname{Phase~Redundancy}
\expandafter\def\csname bibkey@lasser2026micro\endcsname{Microfoundation}
\makeatother

\newtheorem{definition}{Definition}
\newtheorem{lemma}{Lemma}
\newtheorem{theorem}{Theorem}
\newtheorem{remark}{Remark}
\newtheorem{assumption}{Assumption}

\newcommand{\Aadv}{A_{\mathrm{adv}}}
\newcommand{\deltaadv}{\delta_{\mathrm{adv}}}
\newcommand{\deltastar}{\delta_{*}}
\newcommand{\rhogap}{\rho_{\mathrm{gap}}}
\newcommand{\Pproxy}{P}
\newcommand{\Ttruth}{T}
\newcommand{\epsgap}{\varepsilon_{\mathrm{gap}}}
\newcommand{\epsnonres}{\varepsilon_{\mathrm{gap}}^{\mathrm{nonres}}}
\newcommand{\epsfloor}{\varepsilon_{\mathrm{floor}}^{\mathrm{res}}}

\newcommand{\Tbeta}{T_{\beta}}
\newcommand{\Ncasc}{N_{\mathrm{cascade}}}
\newcommand{\Nev}{N_{\mathrm{events}}}
\newcommand{\bclk}{\beta_{\mathrm{clk}}}

\newcommand{\bcal}{\beta_{\mathrm{cal}}}

\title{Lemma 6 ($h_{\mathrm{detect}}$ intensivity): Full Proof Draft v5}
\author{Paper 10 working draft for codex re-review}
\date{}

\begin{document}

\maketitle

\section{Goal and changes from v4}

V4 corrected v3's $\kappa$-scaling error and introduced (C5.HOEFF)
+ strengthened (C11.CLK).  Round~D returned ``v5 needed'': the
$1/\deltastar^2$ scaling is conceptually correct, but v4's
specific formula
$\Tbeta = \max\{2A/\deltastar, R^2\log(1/\beta)/(2\deltastar^2)\}$
is not a rigorous inversion of Lemma~4's exact Hoeffding bound:
substituting $t = R^2\log(1/\beta)/(2\deltastar^2)$ does not give
exponent $\geq \log(1/\beta)$ when $A > 0$.

\textbf{Round D significant problem (Q1, factor-of-4 error).}
For $T \geq 2A/\deltastar$, we have $T\deltastar - A \geq
T\deltastar/2$, hence $(T\deltastar - A)^2 \geq T^2\deltastar^2/4$.
Substituting into the Hoeffding bound:
$2(T\deltastar - A)^2/(TR^2) \geq T\deltastar^2/(2R^2)$.
For this to exceed $\log(1/\beta)$, we need $T \geq 2R^2\log(1/\beta)/
\deltastar^2$ — \emph{four times} v4's coefficient.

V5 uses the conservative form
\[
  \Tbeta(\beta, \deltastar) \;:=\; \max\!\left\{
  \frac{2A}{\deltastar},\;
  \frac{2 R^2 \log(1/\beta)}{\deltastar^2}\right\}
\]
which rigorously gives $\Pr[T_{\mathrm{detect}} > \Tbeta] \leq
\beta$.

\textbf{Round D other issues.}
\begin{itemize}
\item Q2: (C5.HOEFF) ``family-wise corrected'' clause too weak —
  Bonferroni controls error probability, not LLR range. Need
  fixed-cardinality monitored partitions or LLR clipping.
\item Q3: (C11.CLK) sound if $\Ncasc \geq \Tbeta$ is hard audit
  constraint; if empirically estimated, calibration uncertainty
  needs separate $\bcal$ failure probability.
\item Q4: V4 internal labels clean; integration must rewrite
  Lemma~5d and Theorem~1's Layer~2 proof to drop legacy $\beta'$.
\item Q5: Still missing — $\alpha$ explicit clause if multiplicity
  varies; channel multiplicity (C5.MULT); martingale-difference /
  i.i.d. increment structure (Hoeffding requires more than bounded
  range); LLR finiteness/support overlap.
\end{itemize}

\textbf{Changes from v4 (Round D fixes):}

\begin{itemize}
\item \textbf{Q1 fix (rigorous Hoeffding inversion).}  Replaced
  $R^2\log(1/\beta)/(2\deltastar^2)$ with
  $2R^2\log(1/\beta)/\deltastar^2$ (factor of~4).  Step~5 of the
  proof now uses the rigorous chain $T \geq 2A/\deltastar
  \Rightarrow (T\delta - A)^2 \geq T^2\delta^2/4
  \Rightarrow$ exponent $\geq T\delta^2/(2R^2)$, eliminating the
  asymptotic-approximation gap in v4.
\item \textbf{Q2 fix.  (C5.HOEFF) strengthened, (C5.MULT) added.}
  Replaced ``family-wise corrected'' with explicit fixed-cardinality
  monitored partition + uniform LLR clip.  Channel-multiplicity
  policy promoted to its own clause (C5.MULT) for clarity.
\item \textbf{Q3 fix.  Calibration-uncertainty clause $\bcal$.}
  Added explicit treatment: $\Ncasc \geq \Tbeta$ is certified at
  audit time as a hard constraint; calibration error is folded
  into a separate $\bcal$ probability if the operator chooses
  empirical estimation.  Total Layer~2 failure becomes
  $\beta + \bclk + \bcal$ in the most general form, $\beta + \bclk$
  if calibration is certified.
\item \textbf{Q4 fix.  Integration prerequisite list.}  Spelled out
  the Lemma~5d / main\_theorem.tex / deployment\_tooling.tex
  rewrites required at integration: drop $\beta'$ everywhere, use
  $\beta + \bclk + \bcal$ for total Layer~2 failure.
\item \textbf{Q5 fix.  Additional clauses.}
  \begin{itemize}
  \item (C5.IID): SPRT log-likelihood-ratio increments are
    martingale-differences under $p_1$ (or i.i.d. under fixed
    distribution), with bounded conditional variance proxy.
  \item (C5.SUPP): per-channel likelihood ratio satisfies
    $p_0(x), p_1(x) \geq \epsilon > 0$ for all $x$ in the
    monitored partition (or LLR is clipped to $[-R, R]$).
  \item $\alpha$ explicit: monitor Type-I rate, deployment-class
    if (C5.MULT) bounds multiplicity.
  \end{itemize}
\item \textbf{Lemma 4 prose nit.}  Round~D flagged that
  appendix's prose ``$\Lambda_t < A$ is sufficient'' should read
  ``necessary''.  This is an integration item for §A.4, not a
  Lemma~6 issue.
\end{itemize}

\textbf{Carried forward from v4 (unchanged in v5):}
\begin{itemize}
\item Supremum-form $h_{\mathrm{detect}}(\theta) := \sup_s
  \epsgap(s)$.
\item Discrete boundary convention $t_0 := t_0^-$.
\item $\rhogap$ rename + (C11) framing.
\item Union-class KL floor $\deltastar$.
\item (C11.CLK) deterministic floor + $\bclk$ structure.
\end{itemize}

\textbf{Round C other issues.}
\begin{itemize}
\item (C11.CLK) too weak: existence of $N_{\min}, q_{\min} > 0$
  doesn't imply $\Tbeta \leq \tau_{\mathrm{meta}}$ with calibrated
  probability.
\item $\beta'$ mislabeled: v3 used $\beta' = \exp(-\kappa
  \tau_{\mathrm{meta}} \deltastar)$ (the SPRT detection-tail form)
  for the cascade-clock event $\{\Tbeta \leq \tau_{\mathrm{meta}}\}$,
  conflating two distinct probabilities.
\item Constant taxonomy missing: $\alpha$ (Type-I rate), $R$ (LLR
  range), channel multiplicity policy, separate $\bclk$ for clock
  failure, separate from SPRT detection failure.
\item Q2 boundary convention sufficient inside Lemma 6 but must
  surface into Theorem~1's Layer~2 proof at integration time.
\end{itemize}

\textbf{Changes from v3 (Round C fixes):}

\begin{itemize}
\item \textbf{Q1 fix (algebra).}  Define $\Tbeta$ directly from
  Lemma~4's Hoeffding form:
  \[
    \Tbeta(\beta, \deltastar) \;:=\;
    \max\!\left\{\frac{2A}{\deltastar},\;
    \frac{R^2 \log(1/\beta)}{2 \deltastar^2}\right\},
  \]
  where $A := \log((1-\beta)/\alpha)$ is Paper~5's SPRT threshold
  and $R$ is the LLR-increment range from (C5.HOEFF).  In the
  asymptotic regime $\Tbeta \gg A/\deltastar$, the second term
  dominates and gives $\Pr[T_{\mathrm{detect}} > \Tbeta] \leq
  \beta$.  Critical: $\Tbeta$ scales as $1/\deltastar^2$, not
  $1/\deltastar$.
\item \textbf{(C5.HOEFF) added.}  Bounded LLR-increment range
  clause: per-step SPRT log-likelihood-ratio increments are bounded
  in absolute value by $R$, with $R$ deployment-class (not scaling
  with $|P|$ via channel multiplicity).
\item \textbf{(C11.CLK) strengthened.}  Replaced existence-only
  formulation with: there exists deterministic $\Ncasc$ such that
  $\Pr[\Nev(\tau_{\mathrm{meta}}) \geq \Ncasc] \geq 1 - \bclk$, and
  $\Ncasc \geq \Tbeta$ as an audit-time calibration constraint.
  This makes the bound $\Tbeta \leq \tau_{\mathrm{meta}}$
  operational with named failure probability $\bclk$.
\item \textbf{$\beta'$ relabeled.}  V3's use of $\beta' = \exp(-\kappa
  \tau_{\mathrm{meta}} \deltastar)$ is dropped from this lemma's
  statement.  The cascade-clock event $\{\Tbeta \leq
  \tau_{\mathrm{meta}}\}$ now has its own probability bound via
  $\bclk$ from (C11.CLK).  Total Layer~2 failure probability is
  $\beta + \bclk$.
\item \textbf{Constant taxonomy completed.}  $\alpha$ surfaced as
  monitor-design parameter; $R$ surfaced via (C5.HOEFF); $\bclk$
  surfaced via (C11.CLK); channel multiplicity flagged as
  deployment-class via (C5.HOEFF) range constancy.
\item \textbf{Q2 integration note.}  V3's $t_0 := t_0^-$ convention
  carried forward; integration must update Theorem~1's Layer~2
  proof to express the Layer~1$\to$Layer~2 transition in SPRT
  exposure clock.
\end{itemize}

\textbf{Carried forward from v3 (unchanged in v4):}
\begin{itemize}
\item Supremum-form $h_{\mathrm{detect}}(\theta) := \sup_s
  \epsgap(s)$ matching Theorem~1's Layer~2 application.
\item Discrete boundary convention $t_0 := t_0^-$.
\item $\rhogap$ rename and reinforcement that $\rhogap$ excludes
  $K_{\mathrm{Lip}}$.
\item Union-class KL floor $\deltastar = \min(\deltaadv,
  \delta_{\mathrm{adv}}^{\mathrm{env}})$.
\item (C11) framed for promotion to theorem-level condition;
  source-citation correction re Paper~9.
\end{itemize}

\section{Setup: the monitoring clock and $h_{\mathrm{detect}}$}

Theorem~\ref{thm:deployment_safety}'s Layer~2 covers adversarial
strategies $q \in \Aadv \cup \Aadv^{\mathrm{env}}$ during the
detection window
$[t_0, t_0 + T_{\mathrm{detect}}]$.

\paragraph{Boundary convention.}  We define $t_0$ as the
\emph{last SPRT exposure step} at which the deployment is still
within Layer~1's safe region; equivalently, $t_0$ is the SPRT step
immediately preceding the first crossing event ($t_0 := t_0^-$ in
the v2 notation).  $T_{\mathrm{detect}}$ is the number of SPRT
steps from $t_0$ until the SPRT log-likelihood ratio crosses the
detection threshold (Lemma~\ref{lem:4}); the crossing step itself
is counted within $T_{\mathrm{detect}}$.

Under this convention, at $s = s_0 = t_0$ the Layer~1 bound applies
exactly: $\epsgap(t_0) \leq h_{\mathrm{static}}(\theta)$.  The
single-step jump that first crosses Layer~1 occurs at $s_1$, and is
bounded by $\rhogap(\deltastar)$ via (C11).  No ``absorbed buffer''
is required.

\paragraph{Detection-window bound.}
Layer~2 applies $h_{\mathrm{detect}}$ as a point-in-time bound:
\[
  |g(\Ttruth) - g(\Pproxy)|(s) \;\leq\; \mathrm{Lip}(g) \cdot
  h_{\mathrm{detect}}(\theta)
  \quad \text{for all } s \in [t_0, t_0 + T_{\mathrm{detect}}].
\]
Therefore $h_{\mathrm{detect}}$ must bound $\epsgap(s)$ at every
$s$ in the window:
\[
  h_{\mathrm{detect}}(\theta) \;:=\; \sup_{s \in
  [t_0, t_0 + T_{\mathrm{detect}}]} \epsgap(s).
\]

\paragraph{The monitoring clock and (C11.CLK) clock comparability.}
The SPRT machinery (Paper~5) operates on a per-event clock: each
ledger entry triggers an SPRT update, and the per-step KL divergence
drives detection.  All times in this lemma — $t_0$,
$T_{\mathrm{detect}}$, $\Tbeta$, and the cascade-time references
$\tau_{\mathrm{meta}}$, $T_{\mathrm{cascade}}$ — are expressed in
this SPRT exposure clock.

\begin{assumption}[(C5.HOEFF): Bounded LLR-increment range with explicit clipping]
\label{ass:C5HOEFF}
The per-step SPRT log-likelihood-ratio increments
$\ell_n := \log(p_1(X_n)/p_0(X_n))$ are bounded in absolute value
by a deployment-class constant $R$:
\[
  |\ell_n| \;\leq\; R \quad \text{almost surely, for all } n,
\]
with $R$ independent of $|P|$.  The bound is enforced by uniform
LLR clipping: $\ell_n^{\mathrm{used}} := \mathrm{clip}(\ell_n, -R, R)$.
Clipping with deployment-class $R$ ensures the Hoeffding step
applies regardless of underlying-distribution support tails.
\end{assumption}

\begin{assumption}[(C5.MULT): Bounded channel multiplicity]
\label{ass:C5MULT}
The deployment's channel-monitoring policy uses a fixed-cardinality
monitored partition: the number of monitored channels $K_{\mathrm{ch}}$
is a deployment-class constant independent of $|P|$.  In particular,
the four-channel agent-side observable (Paper~5) plus the
environment-side mirror (Paper~10 Lemma~5e) gives
$K_{\mathrm{ch}} \leq 8$; multinomial channel partitions, if used,
must use bounded-cardinality coarsening.
\end{assumption}

\begin{assumption}[(C5.IID): Adapted/i.i.d.\ LLR increments]
\label{ass:C5IID}
Under the alternative hypothesis $p_1$, the per-step LLR increments
$\ell_n$ are either (a) i.i.d.\ with finite mean and bounded range
$R$ ((C5.HOEFF)), or (b) form an adapted martingale-difference
sequence with bounded increments and conditional drift
$\mathbb{E}[\ell_n \mid \mathcal{F}_{n-1}] \geq \delta_n \geq
\deltastar$.  The Hoeffding-Azuma machinery in Lemma~4 applies
under either form.
\end{assumption}

\begin{assumption}[(C5.SUPP): LLR finiteness / support overlap]
\label{ass:C5SUPP}
Either (a) the channel-likelihood pair $(p_0, p_1)$ has overlapping
support with $p_0(x), p_1(x) \geq \epsilon > 0$ for all $x$ in the
monitored partition (with $\epsilon$ a deployment-class constant), or
(b) (C5.HOEFF)'s clipping handles non-overlap by capping the LLR.
This ensures $\ell_n$ is almost-surely finite.
\end{assumption}

\begin{assumption}[(C11.CLK): Clock comparability with calibrated probability]
\label{ass:C11CLK}
There exists a deterministic deployment-class constant $\Ncasc$ and
a calibrated clock-failure probability $\bclk \in (0,1)$ such that:
\[
  \Pr\bigl[\Nev(\tau_{\mathrm{meta}}) \geq \Ncasc\bigr] \;\geq\;
  1 - \bclk,
\]
where $\Nev(\cdot)$ is the SPRT exposure event count up to wall-clock
time $\cdot$.  Furthermore, the operator certifies (at audit time, as
a \emph{hard} constraint) that
\[
  \Ncasc \;\geq\; \Tbeta(\beta, \deltastar)
\]
holds, where $\Tbeta$ is the SPRT detection quantile defined below.

\textbf{Optional empirical-calibration variant.}  If $\Ncasc$ or
$\Tbeta$ are estimated empirically rather than certified
conservatively, the operator must specify a calibration-uncertainty
failure probability $\bcal \in (0,1)$ such that
$\Pr[\Ncasc^{\mathrm{est}} \geq \Ncasc^{\mathrm{true}}] \geq 1 - \bcal$,
and total Layer~2 failure inflates accordingly to $\beta + \bclk +
\bcal$.  In the certified-conservative case, $\bcal = 0$.
\end{assumption}

\textbf{Why (C11.CLK) is strengthened in v5.}  V3's existence-only
form ($N_{\min}, q_{\min}$) did not chain to the
$\Tbeta \leq \tau_{\mathrm{meta}}$ inequality.  V4's deterministic
floor with $\bclk$ chains correctly but treated audit calibration
as implicit.  V5 makes the audit-vs-empirical distinction explicit:
$\bcal = 0$ in the certified case, $\bcal > 0$ when estimating.

\textbf{Why (C5.HOEFF), (C5.MULT), (C5.IID), (C5.SUPP) are
required.}  Lemma~4's Hoeffding step requires four ingredients:
bounded range (C5.HOEFF), constant cardinality (C5.MULT),
adapted/i.i.d.\ structure (C5.IID), and almost-sure finiteness
(C5.SUPP).  V4 conflated some of these; v5 separates them so each
audit hook is targetable.  Channel multiplicity scaling with $|P|$
(broken (C5.MULT)) would inflate $R$ via per-channel cells; absence
of clipping or support overlap (broken (C5.HOEFF) or (C5.SUPP))
would give unbounded LLR increments; non-adapted increments (broken
(C5.IID)) would invalidate Hoeffding-Azuma.

\section{Bounded gap-growth-rate condition (C11)}

\begin{assumption}[(C11): Bounded gap-growth rate per SPRT step]
\label{ass:C11}
For any adversarial strategy $q \in \Aadv \cup
\Aadv^{\mathrm{env}}$ producing channel deviations bounded by
$\deltastar = \min(\deltaadv, \delta_{\mathrm{adv}}^{\mathrm{env}})$
(the Lemmas~5b and~5e KL floors), the per-SPRT-step proxy-truth
gap growth is bounded:
\[
  (\Delta \epsgap)_n \;:=\; \epsgap(s_n) - \epsgap(s_{n-1}) \;\leq\;
  \rhogap(\deltastar),
\]
where $\rhogap(\deltastar)$ is a deployment-class constant
depending on $\mathcal{D}$'s operational parameters (channel
sensitivity-to-gap mapping, event-rate normalization) but
independent of $|P|$.

\textbf{Continuous-time variant.}  Equivalently, in continuous
SPRT exposure clock $s$:
\[
  \frac{d \epsgap}{ds} \;\leq\; \rhogap(\deltastar).
\]
The discrete and continuous forms agree under fixed step duration
or bounded event-rate normalization.
\end{assumption}

\textbf{Operational interpretation.}  Channel deviations bounded by
$\deltastar$ (else SPRT fires faster) limit how much $\epsgap$ can
grow per SPRT step.  The constant $\rhogap$ characterizes the
deployment's channel-deviation-to-gap-growth coupling.

\textbf{$\rhogap$ is calibrated in the $\epsgap$ metric only.}
$\rhogap$ is a bound on the per-step proxy-truth gap process
$\epsgap = $ (Microfoundation's relative sup-norm), \emph{before}
applying any alignment-property mapping $g$.  In particular,
$\rhogap$ excludes the Lipschitz constant $K_{\mathrm{Lip}}$ of $g$;
the Layer~2 operational form multiplies $h_{\mathrm{detect}}$ by
$K_{\mathrm{Lip}}$ separately via (C6).  This eliminates any
double-counting between (C11) gap-dynamics and (C6) alignment-side
smoothness.

\textbf{(C11) is new content.}  \citep[Microfoundation]{lasser2026micro}'s
``$T$ failure modes'' remark notes adequacy/failure modes of the
operational truth $T$ but does not establish a formal per-step
gap-growth-rate bound.  (C11) is a Paper~10 operational assumption,
analogous to (C7) bounded co-evolution and SA1 HHI surrogate
adequacy: deployment-class condition with empirical testability,
not a derived consequence of source-paper machinery.

\section{Statement of Lemma 6 (v5)}

\begin{lemma}[$h_{\mathrm{detect}}$ intensivity]
\label{lem:6}
Under Theorem~\ref{thm:deployment_safety}'s conditions (C5)
continuous SPRT monitoring, (C5.HOEFF) bounded LLR-increment range,
(C5.MULT) bounded channel multiplicity, (C5.IID) adapted/i.i.d.\
LLR increments, (C5.SUPP) LLR finiteness, (C11) bounded gap-growth
rate, and (C11.CLK) clock comparability with calibrated
probability, combined with Lemma~\ref{lem:4} (SPRT Wald--Hoeffding
tail bound) and Lemma~\ref{lem:1} (Layer~1 intensivity):

\textbf{Detection quantile.}  Define the SPRT high-probability
detection quantile via the conservative inversion of Lemma~4's
Hoeffding bound:
\[
  \Tbeta(\beta, \deltastar) \;:=\;
  \max\!\left\{\frac{2A}{\deltastar},\;
  \frac{2 R^2 \log(1/\beta)}{\deltastar^2}\right\},
\]
where $A := \log((1-\beta)/\alpha)$ is Paper~5's SPRT threshold,
$\alpha$ is the Type-I error rate (deployment-class monitor
parameter under (C5.MULT)), $R$ is the LLR-increment range bound
under (C5.HOEFF), $\beta$ is the target detection-tail failure
probability, and $\deltastar = \min(\deltaadv,
\delta_{\mathrm{adv}}^{\mathrm{env}})$ is the union-class KL floor
(Lemmas~\ref{lem:5b}, \ref{lem:5e}).

(i) [Tail bound on $T_{\mathrm{detect}}$.]  $\Tbeta$ rigorously
satisfies $\Pr[T_{\mathrm{detect}} > \Tbeta] \leq \beta$ via
Lemma~\ref{lem:4}'s Hoeffding form (proved in Step~5).

(ii) [Detection-window supremum.]  On the event
$\{T_{\mathrm{detect}} \leq \Tbeta\}$ (probability at least
$1-\beta$):
\[
  h_{\mathrm{detect}}(\theta) \;:=\;
  \sup_{s \in [t_0, t_0 + T_{\mathrm{detect}}]} \epsgap(s)
  \;\leq\; h_{\mathrm{static}}(\theta) + \rhogap(\deltastar) \cdot
  \Tbeta.
\]

(iii) [Lead-time-before-cascade.]  By (C11.CLK)'s deterministic
floor $\Ncasc \geq \Tbeta$ and clock-failure probability $\bclk$
(plus calibration-uncertainty $\bcal$ if applicable):
\[
  \Pr\bigl[\Tbeta \leq \Nev(\tau_{\mathrm{meta}})\bigr] \;\geq\; 1 -
  \bclk - \bcal,
\]
i.e., $\Tbeta \leq \tau_{\mathrm{meta}}$ in SPRT exposure clock
with probability at least $1 - \bclk - \bcal$.  In the certified
case ($\bcal = 0$), this reduces to $1 - \bclk$.

(iv) [Layer~2 total failure bound.]  Total Layer~2 failure
probability (detection failure or cascade-clock miscalibration):
\[
  \Pr[\text{Layer~2 fails}] \;\leq\; \beta + \bclk + \bcal.
\]

\textbf{Intensivity.}  $h_{\mathrm{detect}}(\theta)$ is intensive
in $|P|$ over the deployment class $\mathcal{D}$:
\begin{itemize}
\item $h_{\mathrm{static}}(\theta)$: intensive by Lemma~\ref{lem:1}.
\item $\rhogap(\deltastar)$: intensive by (C11).
\item $\Tbeta = \max\{2A/\deltastar, 2R^2\log(1/\beta)/\deltastar^2\}$:
  intensive because $A$ depends on $\alpha, \beta$ (monitor-design
  parameters; $\alpha$ deployment-class via (C5.MULT)), $R$ is
  deployment-class via (C5.HOEFF), and $\deltastar$ is deployment-
  class via Lemmas~5b/5e.
\item Total bound $h_{\mathrm{static}} + \rhogap \cdot \Tbeta$ is
  a sum of intensive quantities.
\end{itemize}
\end{lemma}

\section{Proof}

\begin{proof}
We bound $h_{\mathrm{detect}}$ as the supremum of $\epsgap$ over
the detection window, derive $\Tbeta$ directly from Lemma~4's
Hoeffding form, and apply (C11.CLK) to bound the cascade-clock
event.

\emph{Step 1 (per-step decomposition).}  Under the boundary
convention from §2, set $s_0 := t_0$ (the last safe SPRT step) and
$s_n := t_0 + n$ in SPRT exposure clock.  For any
$0 \leq n \leq T_{\mathrm{detect}}$:
\[
  \epsgap(s_n) \;=\; \epsgap(s_0) + \sum_{i=1}^{n}
  (\Delta \epsgap)_i.
\]

\emph{Step 2 (apply (C11)).}  By Assumption~\ref{ass:C11}, each
per-step increment satisfies $(\Delta \epsgap)_n \leq
\rhogap(\deltastar)$.  Therefore:
\[
  \epsgap(s_n) \;\leq\; \epsgap(s_0) + n \cdot
  \rhogap(\deltastar).
\]

\emph{Step 3 (boundary condition).}  By the boundary convention,
$s_0 = t_0$ is the last SPRT step at which the deployment is still
within Layer~1's safe region.  Layer~1's intensivity result
(Lemma~\ref{lem:1}) gives:
\[
  \epsgap(s_0) \;=\; \epsgap(t_0) \;\leq\;
  h_{\mathrm{static}}(\theta).
\]
The first crossing event occurs at $s_1$, where (C11) bounds the
single-step jump by $\rhogap(\deltastar)$.

\emph{Step 4 (supremum over window).}  Combining Steps~1--3:
\[
  h_{\mathrm{detect}}(\theta) \;=\;
  \sup_{0 \leq n \leq T_{\mathrm{detect}}} \epsgap(s_n)
  \;\leq\; h_{\mathrm{static}}(\theta) +
  T_{\mathrm{detect}} \cdot \rhogap(\deltastar).
\]

\emph{Step 5 (rigorous Hoeffding inversion to derive $\Tbeta$).}
By Lemma~\ref{lem:4} (\texttt{appendices/proofs.tex} §A.4), the
SPRT detection time satisfies the Hoeffding form:
\[
  \Pr[T_{\mathrm{detect}} > t] \;\leq\;
  \exp\!\left(-\frac{2(t \deltastar - A)^2}{t R^2}\right)
  \quad \text{for } t > A/\deltastar.
\]
We derive $\Tbeta$ such that this bound is $\leq \beta$, using a
rigorous (not asymptotic) chain.

\emph{Substep 5a (regime condition).}  Suppose
$T \geq 2A/\deltastar$.  Then $T\deltastar - A \geq T\deltastar/2$,
so $T\deltastar - A > 0$ and the Hoeffding bound applies.

\emph{Substep 5b (squared lower bound).}  Squaring the inequality
from~5a:
\[
  (T\deltastar - A)^2 \;\geq\; \frac{T^2 \deltastar^2}{4}.
\]

\emph{Substep 5c (exponent lower bound).}  Substituting into the
Hoeffding exponent:
\[
  \frac{2(T\deltastar - A)^2}{T R^2} \;\geq\;
  \frac{2 \cdot T^2 \deltastar^2 / 4}{T R^2} \;=\;
  \frac{T \deltastar^2}{2 R^2}.
\]

\emph{Substep 5d (solve for $T$ such that exponent $\geq \log(1/\beta)$).}
We want
$\frac{T \deltastar^2}{2 R^2} \;\geq\; \log(1/\beta)$,
which gives
\[
  T \;\geq\; \frac{2 R^2 \log(1/\beta)}{\deltastar^2}.
\]

\emph{Substep 5e (conservative $\Tbeta$).}  Combining the regime
condition (Substep~5a) and the tail condition (Substep~5d):
\[
  \Tbeta \;:=\; \max\!\left\{\frac{2A}{\deltastar},\;
  \frac{2R^2 \log(1/\beta)}{\deltastar^2}\right\}.
\]
For any $T \geq \Tbeta$, both the regime condition and the
exponent-magnitude condition hold, so $\Pr[T_{\mathrm{detect}} > T]
\leq \exp(-T\deltastar^2/(2R^2)) \leq \beta$.  In particular,
$\Pr[T_{\mathrm{detect}} > \Tbeta] \leq \beta$.

\emph{Critical: factor of 4 vs.\ v4.}  V4 used $R^2\log(1/\beta)/
(2\deltastar^2)$ from a non-rigorous asymptotic approximation that
implicitly set $A = 0$.  The rigorous chain (Substeps 5a--5d) shows
the correct conservative coefficient is $2R^2/\deltastar^2$, four
times larger.  This is the v5 fix.

\emph{Step 6 (detection-window supremum bound).}  On the event
$\{T_{\mathrm{detect}} \leq \Tbeta\}$ (probability at least
$1-\beta$):
\[
  h_{\mathrm{detect}}(\theta) \;\leq\;
  h_{\mathrm{static}}(\theta) + \rhogap(\deltastar) \cdot
  \Tbeta.
\]
This proves part~(ii) of the lemma.

\emph{Step 7 (cascade-clock event via (C11.CLK)).}  Separately
from the SPRT tail, (C11.CLK) provides:
$\Pr[\Nev(\tau_{\mathrm{meta}}) \geq \Ncasc] \geq 1 - \bclk$, with
hard audit constraint $\Ncasc \geq \Tbeta$ (certified case
$\bcal = 0$) or empirical $\Pr[\Ncasc^{\mathrm{est}} \geq
\Ncasc^{\mathrm{true}}] \geq 1 - \bcal$ (estimated case).
Therefore:
\[
  \Pr[\Tbeta \leq \Nev(\tau_{\mathrm{meta}})] \;\geq\;
  \Pr[\Tbeta \leq \Ncasc \leq \Nev(\tau_{\mathrm{meta}})]
  \;\geq\; 1 - \bclk - \bcal.
\]
In the certified case ($\bcal = 0$), this reduces to $1 - \bclk$.
This proves part~(iii).

\emph{Step 8 (Layer 2 total failure bound).}  By union bound on
the failure modes:
\begin{itemize}
\item Detection-tail $\{T_{\mathrm{detect}} > \Tbeta\}$:
  probability $\leq \beta$ (Step~5).
\item Clock-failure $\{\Tbeta > \Nev(\tau_{\mathrm{meta}})\}$:
  probability $\leq \bclk + \bcal$ (Step~7).
\end{itemize}
\[
  \Pr[\text{Layer~2 fails}] \;\leq\; \beta + \bclk + \bcal.
\]
This proves part~(iv).

\emph{Step 9 (intensivity).}  Each constant is deployment-class
independent of $|P|$:
\begin{itemize}
\item $h_{\mathrm{static}}(\theta)$: intensive by Lemma~\ref{lem:1}
  (under (C6)--(C10)).
\item $\rhogap(\deltastar)$: deployment-class by (C11).
\item $\Tbeta = \max\{2A/\deltastar, 2R^2\log(1/\beta)/\deltastar^2\}$:
  $A = \log((1-\beta)/\alpha)$ depends on $\alpha$ (deployment-class
  via (C5.MULT) bounding multiplicity-correction overhead) and
  $\beta$ (monitor-design parameter); $R$ is deployment-class via
  (C5.HOEFF) clipping; $\deltastar$ is deployment-class via
  Lemmas~5b/5e.  All inputs intensive, so $\Tbeta$ intensive.
\item $\Ncasc, \bclk, \bcal$: deployment-class by (C11.CLK)
  (operator audit-time calibration).
\end{itemize}
Therefore $h_{\mathrm{detect}}(\theta) = h_{\mathrm{static}} +
\rhogap \cdot \Tbeta$ is intensive in $|P|$ over $\mathcal{D}$.
$\Box$
\end{proof}

\section{Discussion}

\paragraph{The two roles of cascade time (carried from v3).}
\begin{itemize}
\item \emph{Integration horizon for $\rhogap$:} $\Tbeta$, the SPRT
  high-probability detection quantile, gives a legitimate
  \emph{upper} bound on $T_{\mathrm{detect}}$ derived from
  Lemma~4's Hoeffding form.
\item \emph{Lead-time-before-cascade guarantee:}
  $\tau_{\mathrm{meta}} \leq T_{\mathrm{cascade}}$ (a \emph{lower}
  bound on cascade time, from
  \citep[Phase~Redundancy]{lasser2026phase}).  In v4, this is
  composed with $\Tbeta$ via (C11.CLK)'s deterministic floor
  $\Ncasc \geq \Tbeta$ and clock-failure probability $\bclk$ —
  \emph{not} via the SPRT detection-tail form (which was v3's
  mistake).
\end{itemize}

\paragraph{Why v3's $\beta'$ labeling was wrong.}
V3 labeled $\beta' = \exp(-\kappa \tau_{\mathrm{meta}} \deltastar)$
as the failure probability for the cascade-clock event $\{\Tbeta
\leq \tau_{\mathrm{meta}}\}$.  But this expression is the SPRT
\emph{detection-tail} form (probability that detection takes longer
than time $\tau_{\mathrm{meta}}$), not a cascade-clock
miscalibration probability.  V4 separates them: $\beta$ is the
detection-tail parameter; $\bclk$ is the clock-calibration failure
parameter (named in (C11.CLK)); their union bound gives total
Layer~2 failure $\leq \beta + \bclk$.

\paragraph{Why v3 and v4 had successive $\Tbeta$ errors.}
\emph{V3 error.}  V3 used $\Tbeta = \log(1/\beta)/(\kappa \delta)$
treating $\kappa$ as $\delta$-independent.  But Lemma~4's
$\kappa = 2\delta/R^2$ contains a factor of $\delta$, so the exponent
$\kappa T \delta = 2T\delta^2/R^2$ scales as $\delta^2$, not
$\delta$.  V4 corrected the scaling.

\emph{V4 error.}  V4 used $\Tbeta = R^2\log(1/\beta)/(2\deltastar^2)$
from a non-rigorous asymptotic approximation.  Substituting this
$T$ into the exact Hoeffding bound gives exponent $\log(1/\beta)
\cdot (1 - A\deltastar/c)^2$ with $c = R^2\log(1/\beta)/2$, strictly
\emph{less than} $\log(1/\beta)$ when $A > 0$.  V5 corrects this by
the rigorous chain in Step~5: $T \geq 2A/\deltastar \Rightarrow
(T\deltastar - A)^2 \geq T^2\deltastar^2/4 \Rightarrow$ exponent
$\geq T\deltastar^2/(2R^2)$.  Solving for exponent $\geq
\log(1/\beta)$ gives $T \geq 2R^2\log(1/\beta)/\deltastar^2$ — four
times v4's coefficient.

\paragraph{Lemma 5d integration item (Round B Q3 follow-up,
carried).}
At integration time, §4.10's Lemma~5d should be restated using the
union-class floor $\deltastar$.  This is a small amendment to §4.10
and §A.4 that should accompany Lemma~6 integration.

\paragraph{Constant taxonomy (v5 complete).}
\begin{itemize}
\item $h_{\mathrm{static}}(\theta)$: intensive composition
  (Lemma~\ref{lem:1}).
\item $\rhogap(\deltastar)$: \emph{Paper 10 operational constant}
  via (C11); calibrated by per-channel sensitivity-to-gap
  measurement (Audit~7).
\item $\Tbeta = \max\{2A/\deltastar, 2R^2\log(1/\beta)/\deltastar^2\}$:
  rigorous SPRT-derived detection quantile.
\item $A = \log((1-\beta)/\alpha)$: Paper~5 SPRT threshold.
\item $\alpha$: Type-I error rate, monitor-design parameter
  (deployment-class via (C5.MULT) bounding multiplicity-correction
  overhead).
\item $\beta$: Type-II error rate / detection-tail parameter,
  monitor-design parameter (deployment-class).
\item $R$: per-step LLR-increment range, surfaced via (C5.HOEFF);
  enforced by clipping $\ell_n \to \mathrm{clip}(\ell_n, -R, R)$.
\item $K_{\mathrm{ch}}$: monitored-channel cardinality, surfaced via
  (C5.MULT) (deployment-class).
\item $\Ncasc$: deterministic event-count floor, surfaced via
  (C11.CLK) (audit-calibrated to satisfy $\Ncasc \geq \Tbeta$).
\item $\bclk$: clock-calibration failure probability ((C11.CLK)).
\item $\bcal$: calibration-uncertainty failure probability
  ((C11.CLK) optional empirical variant); $\bcal = 0$ in certified
  case.
\item $\tau_{\mathrm{meta}}, T_{\mathrm{cascade}}$: source-derived
  cascade-time references from
  \citep[Phase~Redundancy]{lasser2026phase} (lower bound only;
  with the (C7) bounded co-evolution caveat).
\end{itemize}

\paragraph{Layer 2 operational form.}
Combining Lemma~6 with $\mathrm{Lip}(g) \leq K_{\mathrm{Lip}}$ from
(C6):
\[
  |g(\Ttruth) - g(\Pproxy)|(s) \;\leq\; K_{\mathrm{Lip}} \cdot
  \bigl(h_{\mathrm{static}}(\theta) + \rhogap(\deltastar) \cdot
  \Tbeta\bigr)
\]
with probability at least $1 - (\beta + \bclk + \bcal)$ over the
joint event $\{T_{\mathrm{detect}} \leq \Tbeta\} \cap \{\Tbeta \leq
\Nev(\tau_{\mathrm{meta}})\}$.  The right-hand side is the sum of
intensive constants times $K_{\mathrm{Lip}}$, all $|P|$-independent.

\paragraph{Integration prerequisites (Round D Q4 follow-up).}
At appendix integration, the following must accompany the v5 proof:
\begin{itemize}
\item Drop legacy $\beta'$ from §4.10 Lemma~5d statement;
  replace with v5's $\beta + \bclk + \bcal$ structure or restate
  Lemma~5d as a union-class quantile composition.
\item Update Theorem~1 Layer~2 proof in §5
  (\texttt{main\_theorem.tex}) to use $\Tbeta$ instead of
  $T_{\mathrm{cascade}}$ in Step~6, adopt the $t_0 := t_0^-$
  convention, and cite Lemma~6's parts~(ii)--(iv).
\item Drop the v1 source-citation overstatement in
  \texttt{main\_theorem.tex} line~365 (Paper~9's ``$T$ failure
  modes'' remark).
\item Update §A.4 Lemma~4 prose nit: ``$\Lambda_t < A$'' is
  necessary for $T_{\mathrm{detect}} > t$, not sufficient.
\item Add (C5.HOEFF), (C5.MULT), (C5.IID), (C5.SUPP) to
  Theorem~1's condition list (alongside or as sub-clauses of (C5)).
\item Add (C11) bounded gap-growth rate + (C11.CLK) clock
  comparability.
\item Add Audit~7 to §8 deployment tooling: calibrate $\rhogap$,
  $R$, $\Ncasc$, $\bclk$.
\end{itemize}

\section{What this proof does and does not establish}

\textbf{Establishes:}
\begin{itemize}
\item Algebraically correct $\Tbeta$ derivation from Lemma~4's
  Hoeffding form, with explicit $1/\deltastar^2$ scaling.
\item (C5.HOEFF) bounded LLR-increment range as a deployment-class
  condition, ensuring $R$ does not scale with $|P|$ via channel
  multiplicity.
\item (C11.CLK) strengthened with deterministic floor $\Ncasc \geq
  \Tbeta$ and named clock-failure probability $\bclk$, separating
  cascade-clock event from SPRT detection-tail event.
\item Supremum-form $h_{\mathrm{detect}}$ definition matching
  Theorem 1's Layer 2 application.
\item Rigorous discrete-time boundary condition under the
  $t_0 := t_0^-$ convention.
\item Union-class KL floor $\deltastar$ for Layer 2's full
  detection class.
\item Intensive bound on $h_{\mathrm{detect}}$ as $h_{\mathrm{static}}
  + \rhogap \Tbeta$ with high probability $1-\beta$.
\item Layer~2 total failure probability $\leq \beta + \bclk$ via
  union bound.
\item Constant taxonomy with all $|P|$-dependent constants surfaced
  ($\alpha, \beta, R, \Ncasc, \bclk$ all deployment-class via
  named conditions).
\end{itemize}

\textbf{Does not establish (deferred):}
\begin{itemize}
\item Tightness of $\rhogap(\deltastar), R, \Ncasc$.  Operators
  may calibrate tighter bounds for restricted threat models.
\item Existence of deployment configurations satisfying (C5.HOEFF),
  (C11), and (C11.CLK).  Verified empirically (Audit~7).
\item Structural derivation of (C11) from finer source-paper
  machinery.  Currently a Paper 10 operational assumption.
\item Layer~2 proof updates in Theorem~1 (must adopt $t_0 := t_0^-$
  convention and use $\Tbeta$ instead of $T_{\mathrm{cascade}}$ in
  Step~6 of the Layer~2 proof).  This is an integration task.
\end{itemize}

\section{Specific verification questions for v5 codex review}

\begin{enumerate}
\item \textbf{(Q1 rigorous Hoeffding inversion.)}  Step~5
  Substeps~5a--5e gives the rigorous chain $T \geq 2A/\deltastar
  \Rightarrow (T\deltastar - A)^2 \geq T^2\deltastar^2/4 \Rightarrow$
  exponent $\geq T\deltastar^2/(2R^2)$, leading to $\Tbeta =
  \max\{2A/\deltastar, 2R^2\log(1/\beta)/\deltastar^2\}$.  Is this
  chain correct and tight?  In particular:
  \begin{itemize}
  \item Is the squared lower bound $(T\deltastar - A)^2 \geq
    T^2\deltastar^2/4$ from $T \geq 2A/\deltastar$ rigorous (no
    sign issues)?
  \item Is the conservative $\Tbeta$ tight enough operationally,
    or should we use the exact root form
    $\Tbeta = (2A\deltastar + c + \sqrt{c^2 + 4A\deltastar c}) /
    (2\deltastar^2)$ with $c = R^2\log(1/\beta)/2$?
  \item Does the factor of~4 vs.\ v4 affect any downstream bounds
    in unexpected ways?
  \end{itemize}

\item \textbf{((C5.HOEFF), (C5.MULT) split.)}  V5 splits range
  bound from multiplicity policy:
  \begin{itemize}
  \item (C5.HOEFF): $|\ell_n| \leq R$ via clipping, $R$
    deployment-class.
  \item (C5.MULT): $K_{\mathrm{ch}}$ deployment-class with
    fixed-cardinality partition.
  \end{itemize}
  Is the split clean and necessary, or could (C5.MULT) be absorbed
  into (C5.HOEFF) (since clipping handles unbounded LLR even with
  large multiplicity)?  Conversely, do we need additional clauses
  to handle adversarial channel selection (the adversary picking
  which channel to push hardest)?

\item \textbf{((C5.IID), (C5.SUPP) sufficiency.)}  V5 adds:
  \begin{itemize}
  \item (C5.IID): adapted/i.i.d.\ LLR with conditional drift
    $\geq \deltastar$.
  \item (C5.SUPP): $p_0(x), p_1(x) \geq \epsilon > 0$ or LLR
    clipped.
  \end{itemize}
  Are these the right structural assumptions for Hoeffding to
  apply?  In particular, does (C5.IID) need to specify
  conditional-variance proxy bounds (sub-Gaussian or
  sub-exponential)?  Does (C5.SUPP) need an explicit value of
  $\epsilon$ or is the existence sufficient?

\item \textbf{((C11.CLK) certified vs.\ empirical split.)}  V5
  introduces $\bcal$ for the empirical-calibration variant, with
  $\bcal = 0$ in the certified case.  Total Layer~2 failure
  becomes $\beta + \bclk + \bcal$.

  Is the split clean?  In the certified case, the operator audits
  $\Ncasc \geq \Tbeta$ as a hard inequality at calibration time;
  $\Tbeta$ is computed from $\alpha, \beta, R, \deltastar$, all of
  which are deployment-class.  Is this enough, or do we need the
  operator to certify each input constant independently (e.g.,
  certify $\deltastar$ via Lemma~5b/5e empirical floor)?

\item \textbf{(Integration-prerequisite completeness.)}  V5's
  ``Integration prerequisites'' paragraph lists the §4.10
  Lemma~5d, Theorem~1 Layer~2 proof, §A.4 prose nit, condition-list
  amendments, and Audit~7 additions.  Are there other integration
  items we've missed?  In particular:
  \begin{itemize}
  \item Does the §2 notation table need $\Tbeta, \Ncasc, \bclk,
    \bcal, R, K_{\mathrm{ch}}$ surfaced?
  \item Does §3 invariants need updating to reflect (C11.CLK)'s
    event-throughput floor?
  \item Are there §6 cooperative-anchoring implications of the
    Layer~2 failure decomposition?
  \end{itemize}
\end{enumerate}

\end{document}
